\chapter{Detailed Design}

This chapter presents the detailed design of the Intelligent Clinical Document Processing System (ICDPS). Building upon the high-level architecture described in the previous chapter, it examines the internal structure of individual components, the data flow between processing stages, and the design decisions that support system implementation.

The chapter covers dataset organisation, document preprocessing, model architecture, terminology mapping, training procedures, and inference workflows. Collectively, these elements define the technical processes through which clinical documents are transformed into structured information and SNOMED CT coding recommendations.

\section{Dataset and Data Organisation}

The implementation of the ICDPS required data resources supporting document understanding, clinical information extraction, and terminology mapping tasks. Effective organisation of these resources was necessary to support preprocessing, model development, evaluation, and validation activities throughout the project.

The datasets used in this work comprise clinical documents, annotation files, and supporting terminology resources. These datasets were organised to facilitate model training, testing, and performance assessment while maintaining a clear separation between development and evaluation data. This section describes the composition of the datasets, their organisational structure, and the data preparation procedures adopted during system development.

\subsection{Dataset Description}
The primary dataset comprises 40 real-world clinical documents obtained from Easthampstead Surgery, London, United Kingdom. These documents represent the full spectrum encountered in a typical NHS primary care workflow: outpatient correspondence, specialist referral letters, hospital discharge summaries, NHS 111 reports, Accident and Emergency summaries, ambulance service reports, diabetic eye screening reports, and out-of-hours primary care reports. The dataset includes both digitally generated PDF documents and scanned document images produced from photographed or photocopied clinical letters.

Each document contains structured data fields (patient NHS number, date of birth, consultant name, department) and unstructured free-text clinical narrative sections describing presenting complaints, diagnoses, management plans, medication prescriptions, and follow-up recommendations. The dataset is annotated with ground-truth SNOMED CT codes for 25 documents, enabling quantitative evaluation of SNOMED mapping accuracy. Table~\ref{tab:dataset_characteristics} summarises the key dataset characteristics.

\begin{table}[htbp]
\centering
\caption{Dataset Characteristics Summary}
\label{tab:dataset_characteristics}
\begin{tabular}{|p{4.5cm}|p{10cm}|}
\hline
\textbf{Characteristic} & \textbf{Description} \\
\hline
Total Documents & 40 clinical documents \\
\hline
Source Institution & Easthampstead Surgery, London, UK \\
\hline
Document Formats & Native PDF (digital); JPEG/PNG/TIFF (scanned images) \\
\hline
Page Range & 1 to 8 pages per document (mean 3.2 pages) \\
\hline
Annotated Documents & 25 documents with ground-truth SNOMED CT codes \\
\hline
Document Types & 10 categories (discharge, specialist, NHS 111, A\&E, ambulance, etc.) \\
\hline
Noise Characteristics & Skewed scans, inconsistent lighting, varying print quality \\
\hline
Clinical Content & Diagnoses, medications, investigations, referrals, procedures \\
\hline
\end{tabular}
\end{table}

\subsection{Data Collection and Sources}
The dataset was obtained through collaboration with Easthampstead Surgery under an academic research agreement governing data handling, anonymisation, and use within the project scope. All documents were de-identified prior to use, with patient names, NHS numbers, dates of birth, addresses, and other PHI replaced with synthetic identifiers, consistent with NHS Information Governance requirements.

Document collection involved scanning physical paper records at 300 dpi using an NHS-compliant document scanner and exporting digital records from the document management system in PDF format. Additional scanned copies were generated by photographing physical documents under controlled lighting conditions to simulate mobile image capture scenarios commonly encountered in clinical environments. These images incorporated variations in skew, noise levels, contrast, and illumination conditions to reflect the diversity of document quality encountered in real-world clinical environments.

Including documents with diverse image characteristics enabled the preprocessing and OCR components to be assessed under a range of realistic operating conditions rather than being evaluated solely on clean, digitally generated documents.

\subsection{Dataset Structure and Characteristics}

The dataset encompasses documents exhibiting substantial variation in layout, formatting conventions, font styles, and structural organisation. While some records follow predefined templates with clearly identifiable sections and structured data fields, others are composed largely of free-text clinical narratives with limited visual organisation.

Such diversity is characteristic of clinical documentation originating from different healthcare environments and introduces additional complexity for automated information extraction. The presence of heterogeneous document structures influenced the design of the extraction framework, leading to the adoption of a multi-stage processing pipeline rather than a template-dependent approach that would be less adaptable across the range of document formats represented in the dataset.

Content analysis of the 40 documents identified the following clinical entity categories: diagnostic conditions (62\% of documents), medications with dosages (78\%), physiological measurements (45\%), clinical procedures (38\%), anatomical locations (71\%), and temporal references (83\%). The class distribution of SNOMED CT codes in the annotated subset is dominated by disorder/finding concepts (58\%) and drug substance concepts (29\%), with procedure (9\%) and body structure (4\%) concepts forming the minority. Table~\ref{tab:document_type_distribution} presents the distribution of document types.

\begin{table}[htbp]
\centering
\caption{Distribution of Document Types in the ICDPS Dataset}
\label{tab:document_type_distribution}
\begin{tabular}{|p{8cm}|p{2cm}|p{3cm}|}
\hline
\textbf{Document Type} & \textbf{Count} & \textbf{Percentage (\%)} \\
\hline
Hospital Discharge Summaries & 8 & 20.0 \\
\hline
Specialist Clinical Letters & 9 & 22.5 \\
\hline
NHS 111 Reports & 5 & 12.5 \\
\hline
Accident and Emergency Reports & 4 & 10.0 \\
\hline
Ambulance Service Reports & 4 & 10.0 \\
\hline
Private Specialist Letters & 3 & 7.5 \\
\hline
Diabetic Eye Screening Reports & 3 & 7.5 \\
\hline
Out-of-Hours Primary Care Reports & 2 & 5.0 \\
\hline
External Provider Reports (Eye/ENT) & 2 & 5.0 \\
\hline
\textbf{Total} & \textbf{40} & \textbf{100.0} \\
\hline
\end{tabular}
\end{table}

\section{Preprocessing Pipeline Design}

Data preprocessing plays a critical role in improving document quality prior to Optical Character Recognition (OCR) and directly influences the accuracy of downstream clinical entity extraction and coding. Clinical documents frequently contain variations in scan quality, skewed pages, uneven lighting, handwritten notes, stamps, and background artefacts that reduce OCR performance. Therefore, a structured preprocessing pipeline was implemented to improve document readability and ensure consistent input quality before processing with Amazon Textract.

\subsection{Data Cleaning}

Collected clinical documents were reviewed and preprocessed before OCR execution. Blank or nearly empty pages were removed to avoid unnecessary processing overhead. Documents were converted to a standardized RGB image format to ensure compatibility with OpenCV image operations and Amazon Textract input requirements. Image dimensions and resolution were also normalized to reduce variations introduced by different scanning devices and document sources.

Pages containing severe distortions, incomplete scans, or excessive artefacts were identified during preprocessing. Documents with image quality issues likely to cause substantial OCR errors were excluded from quantitative evaluation to prevent unreliable results from affecting performance measurements.

\subsection{Noise Reduction and Handling Missing Information}

Clinical documents often contain handwritten annotations, photocopy artefacts, ink stamps, variations in paper quality, and scanning-related noise that can affect document readability and downstream text extraction processes. These factors can affect OCR accuracy and reduce the quality of downstream information extraction. To reduce their impact, image enhancement techniques implemented with OpenCV were applied before text extraction.

Morphological filtering operations were used to suppress isolated noise while preserving character structures and connected text components. Adaptive Gaussian thresholding was applied to improve text visibility under varying lighting and scanning conditions. These operations increased text contrast and improved OCR readability.

When specific information could not be extracted reliably, missing values were retained as null fields rather than estimated. Similarly, the summarization component was configured to indicate unavailable information instead of generating unsupported clinical content.

\subsection{Text Transformation and Normalization}

Following noise reduction, image transformations were applied to improve OCR consistency. Documents were converted to grayscale representations and processed using adaptive binarization to increase foreground text visibility. Page skew correction was performed by estimating the dominant text orientation and rotating pages to align text horizontally. Rotation limits were applied to prevent excessive corrections on pages containing limited textual content.

After OCR extraction, text-level normalization operations were performed before NLP processing. Multi-page document outputs were merged into a unified text representation while preserving page boundaries. Repeated headers and footers generated during OCR extraction were removed to reduce redundancy. Unicode characters were normalized into a consistent format to ensure compatibility with downstream components. Basic sentence-boundary corrections were also applied where OCR-induced line breaks interrupted sentence continuity.

These preprocessing operations reduce OCR-related noise and provide more consistent input for entity extraction, terminology mapping, and summarization components.

\subsection{Tokenization and Label Alignment}

Clinical text is tokenized using the BioBERT WordPiece tokenizer with a maximum sequence length of 512 tokens. During model fine-tuning, entity annotations from the curated clinical document dataset are aligned with the generated subword token sequence using a token-label mapping strategy. The first subword token corresponding to an annotated entity receives the appropriate B- or I- label according to the BIO tagging scheme, while subsequent subword tokens belonging to the same word receive an X label indicating continuation.

Tokens assigned the X label are excluded from loss computation during training. This approach allows the model to learn entity boundaries from the original word-level annotations without introducing duplicate supervision signals for subword fragments.

For documents exceeding 512 tokens, a sliding-window strategy is applied using a window size of 512 tokens and a stride of 256 tokens. Duplicate predictions within overlapping regions are resolved by retaining predictions from the window in which the token appears within the non-overlapping section. If the confidence score of an overlapping prediction exceeds the retained prediction by more than 15%, the higher-confidence prediction is selected instead.

\subsection{Data Augmentation}

To increase the representation of Allergy and Dosage entities in the training data, two augmentation techniques were applied:

\begin{enumerate}[(i)]

\item \textbf{Entity Synonym Replacement:} Clinical entities were replaced with equivalent terms obtained from SNOMED CT synonym lists, generating additional training examples while preserving the original annotations.

\item \textbf{Back-Translation:} Sentences containing underrepresented entities were translated into French and then translated back into English using a neural machine translation model. This process generated paraphrased variants while preserving the original clinical meaning.

\end{enumerate}

\section{NER Model Architecture Design}
The Named Entity Recognition (NER) component is responsible for identifying clinically relevant information from unstructured medical text. Clinical documents often contain abbreviations, specialized terminology, incomplete sentences, and context-dependent expressions that make entity extraction challenging. To address these characteristics, the proposed system uses a BioBERT-based architecture that generates contextual representations of clinical text and predicts entity labels at the token level.

The model combines transformer-based contextual embeddings with sequence-level classification mechanisms to identify diagnoses, medications, procedures, allergies, laboratory findings, and other clinical entities. Additional contextual analysis is applied during post-processing to distinguish affirmed, negated, and uncertain clinical statements, improving the quality of extracted information before SNOMED CT mapping and downstream processing.

\subsection{BioBERT Encoder}
The NER encoder is BioBERT-base (cased), comprising 12 transformer encoder layers, 12 attention heads per layer, a hidden dimension of 768, and a total of 110 million parameters. The model processes tokenized input and produces a sequence of contextual embeddings of dimension 768 for each token. The \texttt{[CLS]} token embedding is not used for NER; only the per-token embeddings are passed to the classification head.

\subsection{CRF Classification Head}
A linear classification layer maps each 768-dimensional token embedding to a vector of logits over the label vocabulary (17 labels: B- and I- for each of 8 entity types, plus O). A Conditional Random Field (CRF) layer is applied on top of the linear layer to enforce label sequence constraints. The CRF layer learns a transition score matrix $T$ where $T[i][j]$ represents the learned score of transitioning from label $i$ to label $j$. During inference, the Viterbi algorithm decodes the optimal label sequence.

The CRF layer enforces the following hard constraints:
\begin{itemize}
    \item I-TYPE labels cannot follow O labels without an intervening B-TYPE label.
    \item I-TYPE labels of one entity type cannot follow B-TYPE or I-TYPE labels of a different entity type.
    \item B-TYPE and I-TYPE labels of the same entity type can follow each other without restriction.
\end{itemize}

\subsection{Negation Detection Module}

The negation detection module determines whether an extracted clinical entity is affirmed, negated, or expressed with uncertainty. For each identified entity, a context window of $\pm 5$ tokens surrounding the entity span is extracted from the tokenized document. BioBERT embeddings generated for the context window are passed to a linear classification layer that predicts one of three assertion states: \textit{Affirmed}, \textit{Negated}, or \textit{Uncertain}.

The classifier is initialized using NegEx trigger phrase lists and subsequently adapted using assertion annotations derived from the curated clinical document dataset and NHS-representative records. This combination of rule-based initialization and supervised learning allows the model to recognize both explicit negation triggers and more complex contextual expressions of uncertainty.

\section{SNOMED CT Mapping Module Design}

The SNOMED CT mapping module is responsible for linking extracted clinical entities to standardized terminology concepts. This task is complicated by the variability of clinical language, where the same concept may be represented using synonyms, abbreviations, alternative spellings, or context-specific expressions. Consequently, exact keyword matching alone is often inadequate for reliable terminology assignment. To address this challenge, the mapping framework employs a combination of semantic retrieval and candidate ranking techniques to identify the SNOMED CT concepts most closely aligned with the meaning of the extracted entity.

The proposed framework represents clinical entities and SNOMED CT descriptions using vector embeddings. Candidate concepts are retrieved through similarity-based search and subsequently ranked using contextual information derived from the source document. This approach enables the system to generate confidence-ranked coding suggestions while supporting large-scale terminology collections containing hundreds of thousands of concepts.

\subsection{Concept Embedding Index Construction}
The SNOMED CT concept description index is constructed as follows:
\begin{enumerate}[(i)]
    \item All active concepts in the UK SNOMED CT release are extracted.
    \item For each concept, its preferred term, fully specified name, and all active synonyms are concatenated into a description string.
    \item FastText embeddings are computed for each description string using a character $n$-gram model ($n = 3$--$6$) pre-trained on a biomedical text corpus.
    \item The resulting embedding vectors (dimension 300) are indexed in a FAISS \texttt{IndexFlatIP} (inner product) index for exact search, and a FAISS \texttt{IndexIVFPQ} index for approximate search.
\end{enumerate}

\subsection{Candidate Retrieval and Reranking}
For each affirmed entity, the following two-stage retrieval process is performed:
\begin{itemize}
    \item \textbf{Stage 1 --- Candidate Retrieval:} The entity text is encoded using the FastText model to produce a query vector. The top-50 nearest neighbours (by inner product similarity) are retrieved from the FAISS index, producing 50 candidate SNOMED CT concepts.
    \item \textbf{Stage 2 --- Reranking:} The 50 candidates are reranked using a cross-encoder model --- a fine-tuned BioBERT model that takes as input the concatenation of the entity span (with its surrounding context) and the SNOMED CT concept description, and outputs a relevance score. The top-5 candidates by reranking score are returned as the final suggestions with associated confidence scores.
\end{itemize}

\subsection{Confidence Score Calibration}
Raw reranker scores are calibrated to produce reliable probability estimates using Platt scaling, fitting a logistic regression model on the reranker scores and binary relevance labels from a held-out validation set. Calibrated confidence scores are verified to be within 5\% of empirical accuracy across five confidence deciles on the validation set.

\section{Model Architecture and Experimental Setup}

The ICDPS integrates multiple AI/ML model components, each addressing a distinct processing challenge. This section presents the architecture of each model component and the experimental setup adopted for training and evaluation.

\subsection{Model Selection}
Model selection was guided by three criteria: (i) clinical domain suitability, favouring models pre-trained on biomedical or clinical text; (ii) service integration feasibility, preferring managed AWS services; and (iii) published performance benchmarks on comparable tasks. Table~\ref{tab:model_selection} presents the candidate models considered and the selection rationale for each pipeline component.

\begin{longtable}{
|>{\raggedright\arraybackslash}p{2.7cm}
|>{\raggedright\arraybackslash}p{4.5cm}
|>{\raggedright\arraybackslash}p{7cm}|}

\caption{Model Selection Analysis for ICDPS Components}
\label{tab:model_selection}
\\

\hline

\textbf{Component} &
\textbf{Candidates Considered} &
\textbf{Selected Model and Rationale}
\\

\hline
\endfirsthead

\hline

\textbf{Component} &
\textbf{Candidates Considered} &
\textbf{Selected Model and Rationale}
\\

\hline
\endhead

OCR Engine
&
AWS Textract; Tesseract 5; Google Cloud Vision
&
AWS Textract --- superior table/form extraction; per-word confidence scoring; HIPAA-eligible
\\

\hline

Layout Understanding
&
LayoutLMv3; LayoutXLM; DocFormer
&
LayoutLMv3 --- state-of-the-art on PubLayNet/FUNSD; multi-modal text + layout + image
\\

\hline

Medical NER
&
AWS Comprehend Medical; BioBERT; spaCy scispaCy
&
AWS Comprehend Medical --- managed HIPAA-eligible; native InferSNOMEDCT API
\\

\hline

Clinical LLM
&
AWS Bedrock Claude Sonnet; GPT-4; Llama 3
&
AWS Bedrock Claude Sonnet --- best instruction following; lowest hallucination rate; HIPAA-eligible
\\

\hline

Concept Embedding
&
SapBERT; BioBERT-NER; ClinicalBERT
&
SapBERT --- purpose-designed for medical entity normalisation on SNOMED concepts
\\

\hline

Vector Search
&
FAISS; Pinecone; Milvus
&
FAISS --- open-source; GPU-accelerated; sub-100 ms retrieval
\\

\hline

\end{longtable}

\subsection{Model Architecture Design}

\textbf{LayoutLMv3 Architecture:} LayoutLMv3 is a unified multi-modal pre-trained model for document understanding. The model architecture consists of a shared transformer encoder with 12 layers, 12 attention heads, and a hidden dimension of 768 (LayoutLMv3-base). Text tokens are represented as WordPiece embeddings summed with 1D position embeddings and 2D spatial layout embeddings derived from word bounding box coordinates. Image patches are encoded using a linear projection of 16$\times$16 pixel patches from the document page image. Cross-modal attention between text token representations and image patch representations enables the model to associate text with its visual context on the page. In the ICDPS, LayoutLMv3 is fine-tuned for token classification into five layout-semantic categories: Title, Section Header, Body Text, Data Field, and Table Cell.

\textbf{AWS Comprehend Medical NER:} The \texttt{InferSNOMEDCT()} API operates on a clinical BERT-based sequence labeller pre-trained on MIMIC-III clinical notes, PubMed abstracts, and UMLS metathesaurus concept descriptions. The model applies a CRF output layer to produce entity span labels in BIO format. Each detected entity is associated with a SNOMED CT concept ID, a description string, a clinical category, and a detection confidence score.

\textbf{AWS Bedrock Claude Sonnet:} Claude Sonnet provides a 200,000-token context window, enabling processing of multi-page clinical documents without truncation. Three role-differentiated system prompts are constructed at runtime: a Clinician System Prompt for structured clinical summaries covering findings, diagnoses, and recommended actions; a Patient System Prompt for plain-English explanation avoiding medical jargon; and a Pharmacist System Prompt focusing on medication details, dosages, and contraindications.

\subsection{Training Configuration}
The LayoutLMv3 fine-tuning training configuration was established through a preliminary hyperparameter grid search. Table~\ref{tab:training_config} presents the final training configuration adopted.

\subsection{Experimental Workflow}
The experimental workflow followed structured training, evaluation, and hyperparameter refinement cycles. In each cycle: (i) LayoutLMv3 was fine-tuned on the augmented training dataset; (ii) the fine-tuned model was evaluated on the validation set using token-level accuracy, precision, recall, and F1-score for each zone class; (iii) training and validation loss curves were inspected for overfitting; (iv) hyperparameters were adjusted based on validation performance; and (v) the best checkpoint by macro F1-score was retained. For SNOMED mapping, layer-by-layer ablation experiments quantified the marginal contribution of each fallback layer to overall SNOMED code recall. UCS weighting coefficients were validated through threshold sensitivity analysis across the full confidence threshold range $[0.60, 0.80]$ for all 21 document type categories.

\begin{table}[htbp]
\centering
\caption{LayoutLMv3 Fine-Tuning Training Configuration}
\label{tab:training_config}
\begin{tabular}{|p{5cm}|p{8cm}|}
\hline
\textbf{Hyperparameter} & \textbf{Value} \\
\hline
Base Model & \texttt{microsoft/layoutlmv3-base} (Hugging Face Hub) \\
\hline
Task & Token Classification --- document layout segmentation \\
\hline
Training Epochs & 15 (early stopping patience: 5 epochs) \\
\hline
Batch Size & 8 (gradient accumulated over 4 steps; effective batch = 32) \\
\hline
Optimiser & AdamW \\
\hline
Initial Learning Rate & $2 \times 10^{-5}$ \\
\hline
Learning Rate Schedule & Cosine annealing with warm-up (warmup ratio = 0.1) \\
\hline
Weight Decay & 0.01 \\
\hline
Dropout Rate & 0.1 \\
\hline
Max Input Sequence Length & 512 tokens \\
\hline
Image Resolution & $224 \times 224$ pixels (ViT patch size $16 \times 16$) \\
\hline
Hardware & NVIDIA A10G GPU (AWS g5.xlarge) \\
\hline
Training Duration & Approximately 4.5 hours for 15 epochs \\
\hline
Framework & PyTorch 2.1 + Hugging Face Transformers 4.37 \\
\hline
\end{tabular}
\end{table}

\subsection{Validation Strategy}
The primary validation strategy for SNOMED CT mapping accuracy employs leave-one-out cross-validation (LOOCV) across the 25 annotated documents, reporting precision, recall, and F1-score at the entity-code pair level. LOOCV was selected because the small dataset size makes $k$-fold partitions unstable for minority document categories. For LayoutLMv3 zone classification, stratified 5-fold cross-validation on the synthetic training dataset reports token-level macro F1-score. OCR accuracy is evaluated using character error rate (CER) and word error rate (WER) computed against manually verified ground-truth transcriptions of 15 representative document pages.

\section{Inference and Post-Processing Design}

After model training, the ICDPS processes clinical documents through an inference pipeline that generates entity predictions, terminology mappings, summaries, and structured outputs. In addition to model inference, several post-processing stages are applied to improve output quality and ensure consistency across downstream components.

These stages encompass prediction refinement, confidence assessment, terminology ranking, and the generation of structured outputs suitable for downstream processing and review. Together, they transform raw model predictions into outputs that can be reviewed and used within clinical workflows.

\subsection{Optimisation Techniques}

LayoutLMv3 fine-tuning uses the AdamW optimizer, which separates weight decay from the adaptive gradient update process and is widely used for transformer-based models. A weight decay value of 0.01 is applied to all parameters except bias terms and layer-normalization parameters. Learning-rate warm-up is performed over the first 10% of training steps to reduce instability during the initial stages of training.

Cosine annealing is used to gradually reduce the learning rate throughout training, allowing the model to converge more smoothly. Gradient clipping with a maximum norm of 1.0 limits excessively large gradient updates. Mixed-precision training (FP16) reduces GPU memory consumption by approximately 40% and allows a larger effective batch size to be used during training. Early stopping with a patience value of five validation epochs is applied to reduce overfitting and prevent unnecessary training iterations.

\subsection{Loss Functions and Evaluation Logic}

The LayoutLMv3 token-classification model is trained using a weighted cross-entropy loss function. Class weights are assigned inversely proportional to class frequency in order to reduce the impact of class imbalance between frequently occurring Body Text tokens and less common classes such as Table Cell and Section Header tokens.

The weighted loss function, denoted by $\mathcal{L}_w$, is defined as:
\begin{equation}
\mathcal{L}_w = -\sum_{c} \left[ w_c \cdot y_c \cdot \log(p_c) \right] \quad \text{where} \quad w_c = \frac{N_{\text{total}}}{N_{\text{classes}} \cdot N_c}
\label{eq:weighted_loss}
\end{equation}

The Unified Confidence Score (UCS) is computed using document-type-dependent weighted linear combinations. For standard NHS documents (discharge summaries, specialist letters):
\begin{equation}
\text{UCS} = 0.40 \times \text{textract\_conf} + 0.20 \times \text{snomed\_conf} + 0.40 \times \text{llm\_conf}
\label{eq:ucs_standard}
\end{equation}

For ambulance and NHS 111 reports, where all-caps tables render SNOMED mapping unreliable:
\begin{equation}
\text{UCS} = 0.50 \times \text{textract\_conf} + 0.50 \times \text{llm\_conf}
\label{eq:ucs_ambulance}
\end{equation}

The UCS is compared against document-type-specific thresholds defined in \texttt{config/\allowbreak document\_type\_\allowbreak config.py}, which defines 21 NHS document categories with per-type review thresholds ranging from 0.62 to 0.72.

\subsection{Inference Pipeline}
The production inference pipeline is implemented as a Python class (\texttt{Icdps\-Inference\-Pipeline}) that encapsulates the complete processing workflow. The pipeline accepts a document path as input and executes the following steps in sequence: document format detection, text extraction (direct or OCR), text normalization, sliding-window tokenization, NER inference, entity span decoding, negation detection, role categorization, summarization, SNOMED CT candidate retrieval, and reranking. The pipeline returns a structured JSON object containing the document metadata, extracted entities, summary, and SNOMED CT suggestions.

\subsection{Post-Processing Rules}

The following rules are applied to NER predictions before structured output generation:

\begin{itemize}

\item \textbf{Contiguous Entity Merging:} Adjacent entities of the same type separated only by whitespace or the conjunction `and'' are merged into a single entity span. For example, the sequence `left'' + `ventricular'' + `hypertrophy'' is merged into the entity ``left ventricular hypertrophy''.

\item \textbf{Minimum Entity Length Filter:} Entity spans containing fewer than two non-whitespace characters are removed, as these are typically generated by tokenization or OCR artefacts.

\item \textbf{Duplicate Entity Deduplication:} Repeated occurrences of the same entity text within a document are consolidated. The first occurrence and its associated SNOMED CT mapping are retained in the structured output, while all occurrences remain highlighted within the document view.

\item \textbf{Negated Entity Exclusion:} Entities classified as \textit{Negated} by the assertion detection module are excluded from SNOMED CT code generation. These entities remain visible to the user but are clearly marked and omitted from exported coding reports.

\end{itemize}

\subsection{Output Schema}

The inference process produces a structured JSON output containing extracted entities, SNOMED CT mappings, confidence scores, generated summaries, and associated metadata. The structure of the output is illustrated in the following schema:

\begin{verbatim}
{
  "document_id": string,
  "document_type": string,
  "extraction_timestamp": ISO8601,
  "entities": [{
    "entity_id": string,
    "text": string,
    "category": string,
    "assertion": string,
    "role_action": string,
    "start_char": int,
    "end_char": int,
    "snomed_suggestions": [{
      "concept_id": string,
      "preferred_term": string,
      "confidence": float
    }]
  }],
  "summary": string,
  "processing_metadata": {
    "model_version": string,
    "processing_time_ms": int
  }
}
\end{verbatim}

\section*{Summary}

This chapter presented the detailed design of the Intelligent Clinical Document Processing System (ICDPS), covering dataset organisation, preprocessing workflows, data augmentation methods, model architecture, terminology mapping procedures, and inference workflows. It also described the post-processing mechanisms and confidence-based evaluation strategies used to improve the reliability and consistency of generated outputs.

Key aspects of the design included the BioBERT-based clinical entity extraction framework, the SNOMED CT retrieval and ranking process, and the handling of lengthy clinical documents through sliding-window processing. Training configurations, optimisation approaches, and output-generation mechanisms were also detailed to establish the technical foundation for system implementation.

Chapter 6 focuses on the implementation of the proposed design and discusses the technologies, frameworks, and development practices used during system construction.
